% 自测题1：基础力学选择1.1.1—1.1.10
% 《大学物理学学习指导》第5版，印刷45—46页／PDF53—54页。
\subsection{自测题1：运动学、动量和弹簧选择题}

\sourcetest{1.1.1}{45}{53}
\problemtext 沿$x$轴运动的质点有$x(t)=3t-5t^3+6$（SI），判断加速度方向和变化性质。
\workings
速度$v_x=3-15t^2$，加速度
\[
\boxed{a_x=-30t}.
\]
对$t>0$它指向$x$负方向且大小随时间增大，并非匀加速，故选\textbf{D}。
\crosscheck $t=0$瞬间$a_x=0$；说“向负方向”指题意通常考察的$t>0$区间。

\sourcetest{1.1.2}{45}{53}
\problemtext 曲线运动中用位矢$\bm r$、路程$s$、速度大小$v$及切向加速度大小$a_\tau$判断四个微分式。
\workings
速度的定义为$\bm v=d\bm r/dt$，而标量路程的变化率
$ds/dt=|\bm v|=v$；加速度向量$d\bm v/dt=\bm a$，
但其模通常并非切向加速度大小：
\[
|\bm a|=\sqrt{a_\tau^2+a_n^2}\ne a_\tau.
\]
因此只把$ds/dt=v$视为题面无歧义的标量恒等式，
对应\textbf{D（仅第3项）}。
\crosscheck 原书其余选项将速度矢量、速度标量与曲线运动的切向量混用。
若将第2式明确定义为\emph{矢量}等式$d\bm r/dt=\bm v$，它当然成立；
但此与原题印刷未加粗的$v$记号不同，属于须保留的记号歧义。

\sourcetest{1.1.3}{45}{53}
\problemtext 速度满足$dv/dt=-kv^2t$，$k>0$，初值$v(0)=v_0>0$，求$v(t)$。
\workings
分离变量
$dv/v^2=-kt\,dt$，
积分得
\[
-\frac1v=-\frac{kt^2}{2}-\frac1{v_0},
\qquad
\boxed{\frac1v=\frac{kt^2}{2}+\frac1{v_0}}.
\]
选\textbf{C}。
\crosscheck $t=0$恢复初值；对$t>0$、$v>0$时速度下降。

\sourcetest{1.1.4}{45}{53}
\problemtext 水平粗糙面上的物块受固定大小$F$、向右上方夹角$\theta$的拉力，摩擦系数$\mu$，求最大水平加速度时的$\theta$。
\workings
在物块始终与地面接触的范围内
$N=mg-F\sin\theta\ge0$，
\[
ma_x=F\cos\theta-\mu(mg-F\sin\theta).
\]
对角度求导，
$-F\sin\theta+\mu F\cos\theta=0$，
给
\[
\boxed{\tan\theta=\mu}.
\]
按题面答案选\textbf{C}。
\crosscheck 必须另检验所求角对应$N\ge0$；
若$F$很大使物体离地，摩擦模型不再适用。

\sourcetest{1.1.5}{45}{53}
\problemtext 光滑半球形碗内半径$R=0.10$m，
小球相对碗静止于离碗底$h=0.04$m处，
碗绕竖直轴匀速旋转，求角速度（取$g=9.8\,\mathrm{m/s^2}$）。
\workings
球心比小球高$R-h=0.06$m，
转动半径为$\rho=\sqrt{2Rh-h^2}$。
法向支持力沿球面法线，其竖直平衡与水平向心分量给
\[
N\frac{R-h}{R}=mg,\qquad
N\frac{\rho}{R}=m\omega^2\rho.
\]
消去$N$后
\[
\boxed{\omega=\sqrt{\frac g{R-h}}
=\sqrt{\frac{9.8}{0.06}}
\simeq12.78\,\mathrm{rad/s}}.
\]
选\textbf{A（约13 rad/s）}。
\crosscheck 原书位置高度是离\emph{碗底}4cm，不是离球心4cm。

\sourcetest{1.1.6}{45--46}{53--54}
\problemtext 质量2kg粒子在$\bm F=12t\bm i$（SI）作用下从原点静止开始运动，求$t=3$s时动量。
\workings
动量定理$\bm p(3)=\int_0^3\bm F\,dt$，
\[
\boxed{\bm p=12\int_0^3t\,dt\,\bm i
=54\bm i\,\mathrm{kg\,m/s}}.
\]
选\textbf{B}。
\crosscheck 力对时间积分直接为动量，不需再次乘物体质量。

\sourcetest{1.1.7}{46}{54}
\problemtext 小球在水平面内半径$R$的圆周作逆时针匀速率$v$运动，
从右侧点A（速度向上）走半周至左侧点B（速度向下），求动量增量。
\workings
图中$\bm p_A=mv\bm j$，$\bm p_B=-mv\bm j$，因此
\[
\boxed{\Delta\bm p=\bm p_B-\bm p_A=-2mv\bm j}.
\]
选\textbf{B}。
\crosscheck 动量是矢量；走半周后速率没变，但动量发生大小$2mv$的变化。

\sourcetest{1.1.8}{46}{54}
\problemtext 两个弹簧$k_A,k_B$串联，静止悬挂物块；求两弹簧弹性势能之比。
\workings
轻弹簧串联受相同拉力$T$，伸长量$x_i=T/k_i$，
\[
U_i=\frac12 k_ix_i^2=\frac{T^2}{2k_i}.
\]
故
\[
\boxed{\frac{U_A}{U_B}=\frac{k_B}{k_A}}.
\]
选\textbf{C}。
\crosscheck 同一拉力下较软弹簧储存的能量更大，不能用$k_A/k_B$直接代替。

\sourcetest{1.1.9}{46}{54}
\problemtext 光滑水平面上等质量物体P冲向带弹簧挡板的静止物体Q，碰后P停下、Q以原来P速度离开。何时弹簧压缩量最大？
\workings
压缩到最大时弹簧伸缩速率为零，
即两物体瞬时\emph{速度相等}，
相对动能全部转换成弹簧弹性势能：
\[
\boxed{v_P=v_Q}.
\]
选\textbf{B}。
\crosscheck P最终停下是碰撞完成后的事件，
不能把它等同最大压缩的时刻。

\sourcetest{1.1.10}{46}{54}
\problemtext 质量$M$重物悬于定滑轮右端，质量$m=M/2$的人抓左端绳向上爬，其相对绳的向上加速度为$a_0$，求人对地加速度（向上正）。
\workings
设绳左段加速度为$a_L$，重物向上的加速度为$-a_L$，
张力为$T$。对人和重物：
\[
ma_P=T-mg,\qquad
M(-a_L)=T-Mg,\qquad a_P-a_L=a_0.
\]
联立得
\[
a_P=\frac{M a_0+(M-m)g}{M+m}.
\]
代入$m=M/2$，
\[
\boxed{a_P=\frac{2a_0+g}{3}}.
\]
选\textbf{A}。
\crosscheck 必须明确“相对绳上爬”与“对地加速度”不同；
不能直接写$a_P=a_0$。
